\begin{textsample}{NEG Dim 4 – 1950 – Score -6.00 – 1950/tv\_com\_1950\_108.txt}  \label{ex:f4_neg_005}
The commercial is in black and white and presents a sequence of laboratory, industrial, and automotive testing scenes, ending with a brand logo.
At the beginning, a gloved hand holds a glass beaker or measuring vessel containing a cloudy, foaming, or milky liquid.
The shot is close-up and laboratory-like, with the vessel dominating the frame.
After a dark transitional shot, a man in a suit appears standing near or in front of a door.
The image is soft and slightly blurred ; a small sign on the door is visible but not legible in the sampled frame.
The scene then shifts to industrial testing \textbf{equipment}.
A \textbf{mechanical} apparatus with two large articulated arms stands in front of what looks like a chamber or heavy framed opening.
The camera then moves closer to the \textbf{equipment}, showing cylindrical joints, rods, and a central \textbf{mechanism}.
The setting appears like a testing laboratory or engineering facility.
Next, close-up shots show a hand operating a \textbf{control} panel.
One shot shows a hand turning or pointing at a labeled “ START SWITCH. ” Another shot shows a finger pressing a “ START ” button.
Other visible labels include “ DYN MOTOR, ” “ JACKET WATER, ” and “ LUBE OIL, ” with additional “ START ” and “ STOP ” buttons arranged in rows.
These shots emphasize manual activation of machinery.
The commercial then cuts back to close-ups of \textbf{mechanical} components : pipes, hoses, metal fittings, and a chamber-like device with bolts around an opening.
A gloved or sleeved worker ’s arm is partly visible near the machinery.
A wheel or tire is then shown pressing or rolling over a patterned or textured surface.
The tire has a bright, high-contrast tread pattern, and the surface beneath it also shows raised or dimpled texture.
This suggests a test setup, though the exact purpose is not explicitly shown.
The next scene shows a row of large machine \textbf{units} behind a barrier or rail.
Each \textbf{unit} has circular reels or wheel-like components and \textbf{control} panels.
The \textbf{equipment} appears to be part of a testing or monitoring installation.
The commercial then cuts outdoors to a parking or service area beside a large building.
Several mid-century cars are parked in a line, each near posts and hoses or cables, suggesting a controlled fueling, testing, or service arrangement.
The exact function is not clear from the frames alone.
Inside a \textbf{control} room, a man sits at a large console with multiple \textbf{dials}, gauges, and meters.
He faces away from the camera and appears to be operating or monitoring the \textbf{controls}.
This is followed by a brief superimposed or transitional shot combining the control-room view with another industrial interior containing large pipes or tanks.
The tire-on-textured-surface shot appears again, reinforcing the testing motif.
Near the end, a suited male \textbf{presenter} appears in front of a large oval sign.
The word “ HUMBLE ” is visible in large capital letters behind him, and a smaller oval also contains the word “ HUMBLE. ” A bright starburst or spark-like graphic appears at the upper left of the oval.
In the final frame, the screen resolves to a clean, centered oval logo with the word “ HUMBLE ” in bold capital letters.
A visible “ OsbornTramain ” watermark appears in the lower-right area of many frames.

% matched lemmas: control, dial, equipment, mechanical, mechanism, presenter, unit
\end{textsample}
